\section{21.05.2026 - Kolleg*innen sind auch nur Menschen}


\subsection{Wortschatz}

\begin{vocabtable}
\deword{sich unterhalten mit + dat}{to talk / to have a conversation with}{Er unterhält sich gern mit seinen Kolleginnen und Kollegen.}
\deword{plaudern mit + dat}{to chat}{In der Pause plaudern die Kollegen über den Alltag.}
\deword{schwätzen mit + dat}{to chat / to gossip}{Sie schwätzt oft mit ihren Kolleginnen in der Küche.}
\deword{angenehm}{pleasant / comfortable}{Eine angenehme Atmosphäre ist am Arbeitsplatz wichtig.}
\deword{das Ereignis, die Ereignisse}{event / occurrence}{Positive Ereignisse im Team stärken die Zusammenarbeit.}
\deword{locker}{relaxed / casual}{Eine lockere Unterhaltung kann das Arbeitsklima verbessern.}
\deword{anwenden + akk}{to apply / to use}{Diese Strategie wird oft in Gesprächen angewendet.}
\deword{geteilt}{shared / divided}{Geteilte Erfahrungen können Kolleginnen und Kollegen verbinden.}
\deword{sich austauschen mit + dat über + akk}{to exchange views / to talk with someone about something}{In der Pause tauschen wir uns über aktuelle Projekte aus.}
\end{vocabtable}


\subsection{Text -- Korrektur}

- Eine andere \wrong{moeglichkeiten}
\correct{Möglichkeit} ist\correct{,} dass sie \wrong{ueber}
\correct{über} \wrong{ihre Beirech}
\correct{ihren Bereich} sprechen.

- \wrong{So viel}
\correct{Soweit} ich \wrong{weiss}
\correct{weiß}, \wrong{sie koennen}
\correct{können sie} \wrong{uber}
\correct{über} \wrong{ihre Kindern}
\correct{ihre Kinder} diskutieren.
